\documentclass[10pt]{amsart}
\usepackage{calc,amssymb,amsmath,ulem,stmaryrd,fullpage}
\usepackage{enumerate}
\usepackage{alltt}
\usepackage{multicol}
\RequirePackage[dvipsnames,usenames]{color}
\let\mffont=\sf
\normalem
%\input{kmacros3.sty}
\input{xy}
\xyoption{all}
\usepackage{hyperref}
\usepackage{graphicx}
\usepackage[all,cmtip]{xy}
\usepackage{verbatim}
\usepackage[left=0.95in,top=0.95in,right=0.95in,bottom=0.95in]{geometry}
\newcommand{\tauCohomology}{T}
\newcommand{\FCohomology}{S}


\theoremstyle{theorem}
%\newtheorem*{mainthm*}{Main Theorem}

\newcommand{\note}[1]{\marginpar{\sffamily\tiny #1}}

\def\todo#1{\textcolor{Mahogany}%
{\sffamily\footnotesize\newline{\color{Mahogany}\fbox{\parbox{\textwidth-15pt}{#1}}}\newline}}

\renewcommand{\O}{\mathcal O}
\newcommand{\ie}{{\itshape i.e.} }
\newcommand{\roundup}[1]{\lceil #1 \rceil}
\newcommand{\rounddown}[1]{\lfloor #1 \rfloor}
\renewcommand{\to}[1][]{\xrightarrow{\ #1\ }}
\newcommand{\onto}[1][]{\protect{\xrightarrow{\ #1\ }\hspace{-0.8em}\rightarrow}}
\newcommand{\into}[1][]{\lhook \joinrel \xrightarrow{\ #1\ }}
%\newcommand{DBDual}[1]{{\underline{\omega}_{#1}}}
\newcommand{\DBDual}[1]{{{\uuline \omega}{}^{\mydot}_{#1}}}
\newcommand{\Tt}{{\mathfrak{T}}}
%\newcommand{\bn}{{\mathfrak{n}} }
\newcommand{\sol}[1]{ \vskip 6pt{\color{blue} {\bf Solution:} #1} \vskip 6pt}

\newcommand{\bR}{{\mathbb{R}}}
\newcommand{\va}{{\vec{a}}}
\newcommand{\vb}{{\vec{b}}}
\newcommand{\vzero}{{\vec{0}}}
\newcommand{\vi}{{\vec{i}}}
\newcommand{\vj}{{\vec{j}}}
\newcommand{\vk}{{\vec{k}}}
\newcommand{\vh}{{\vec{h}}}
\newcommand{\vr}{{\vec{r}}}
\newcommand{\vu}{{\vec{u}}}
\newcommand{\vv}{{\vec{v}}}
\newcommand{\vw}{{\vec{w}}}
\newcommand{\vx}{{\vec{x}}}
\newcommand{\vy}{{\vec{y}}}
\newcommand{\vz}{{\vec{z}}}
\newcommand{\vT}{{\vec{T}}}
\newcommand{\vN}{{\vec{N}}}
\newcommand{\vF}{{\vec{F}}}
\newcommand{\vG}{{\vec{G}}}
\newcommand{\bF}{\mathbb{F}}
\newcommand{\bQ}{\mathbb{Q}}
\newcommand{\bZ}{\mathbb{Z}}
\newcommand{\bC}{\mathbb{C}}


\begin{document}
\title{Worksheet \#2-- Math 5320\\Spring 2020}
\author{Due Friday, January 31st, in Gradescope}

\maketitle


\noindent
{\bf 1.}  An \emph{automorphism} of a ring $R$ is a ring isomorphism $\phi : R \to R$.  Let $A$ be a ring and consider the map $\phi : A[x,y] \to A[x,y]$ which sends elements of $A$ to themselves, sends $x \mapsto y$ and $y \mapsto x$.  Prove that $\phi$ is a ring homomorphism and an isomorphism.  
\vskip 3pt
\emph{Hint:}  You may use material from the text.
\newpage
\noindent
{\bf 2.}  Find all ring automorphisms of the polynomial ring $\mathbb{Z}[x]$.
\newpage
\noindent
{\bf 3.}  Let $R$ be a ring of prime characteristic $p$ (in other words, $p = 0$, where $p$ is $1+1+ \dots + 1$ repeated $p$-times, $\bZ/p\bZ$ has characteristic $p > 0$).  Consider the ring map 
\[
    F : R \to R
\]
defined by $F(r) = r^P$.  Prove that $F$ is a ring homomorphism.  It is called the \emph{Frobenius} homomorphism.
\newpage
\noindent
{\bf 4.}  For each of the following, produce an example of a ring of prime characteristic $p > 0$ where:
\begin{itemize}
    \item[(a)]  The Frobenius homomorphism is \emph{bijective}.
    \item[(b)]  The Frobenius homomorphism is \emph{injective} but \emph{not surjective}.
    \item[(c)]  The Frobenius homomorphsim  is \emph{surjective} but \emph{not injective}.
\end{itemize}
\newpage
\noindent
{\bf 5.}  Suppose $S = \mathbb{Q}[x]$ and let $f = x^4 + x^3 + x^2 + x + 1$.  Let $R = S/(f)$ and suppose $\overline{x}$ denotes $x + (f)$ (the image of $x$ in $R$).  Express the element 
\[
    (\overline{x}^3 + \overline{x}^2 + \overline{x})(\overline{x}^5 + \overline{1})
\]
as a $\bQ$-linear combination of $\{\overline{1}, \overline{x}, \overline{x}^2, \overline{x}^3\}$.  
\newpage
\noindent
{\bf 6.}  Is there a field $k$ such that $k[x]/(x^2)$ is isomorphic (as a ring) to $k[x]/(x^2-1)$.
\newpage
\noindent
{\bf 7.}  Let $R$ be a ring and suppose $a \in R$.  Let $R' = R[x]/(ax - 1)$.  Let $\alpha$ denote the image of $a$ in $R'$.  Note that $\alpha$ is an inverse of $x$.
\begin{itemize}
    \item[(a)]  Show that every element in $R'$ can be written in the form $\alpha^k b$ for some $b \in R$.
    \item[(b)] Prove that the kernel of $R \to R'$ (sending an element to its coset) is the set of $b \in B$ such that $a^n b = 0$ for some $n > 0$.
    \item[(c)] Prove that $R'$ is the zero ring if and only if $a$ is nilpotent\footnote{If a power of $a$ is zero in $R$.}.
\end{itemize}
\newpage
\noindent
{\bf 8.}  Suppose we adjoin an element $\alpha$ such that $\alpha^2 = 1$ to the ring $\mathbb{R}$.  Prove that the resulting ring is isomorphic (as a ring) to $\mathbb{R} \times \mathbb{R}$.  
\newpage
\noindent
{\bf 9.}  In each case, describe the ring obtained by adjoining an element $\alpha$ to $\bF_2 \cong \bZ/2\bZ$ satisfying the given relation.
\begin{itemize}
    \item[(a)]  $\alpha^2 + \alpha + 1 = 0$.
    \item[(b)]  $\alpha^2 + 1 = 0$.
    \item[(c)]  $\alpha^2 + \alpha = 0$.
\end{itemize}
In particular, in each case, identiy the zero divisors, the nilpotent elements, and the units.


\end{document} 